\pagenumbering{roman}
\setcounter{page}{1}

\selecthungarian
%----------------------------------------------------------------------------
% Abstract in Hungarian
%----------------------------------------------------------------------------
\chapter*{Kivonat}\addcontentsline{toc}{chapter}{Kivonat}

\noindent Diplomamunkám témája egy komplex, \textbf{Unity} alapú ökoszisztéma-szimulációs keretrendszer fejlesztése, amelynek elsődleges célja a populációgenetikai folyamatok és az öröklődő tulajdonságok vizsgálata digitális környezetben. A fejlesztés mostanra egy négyéves intervallumot ölel fel, amely során a szoftver egy egyszerű ágensalapú modellből egy robusztus, eseményvezérelt rendszerré fejlődött.

A projekt kezdeti szakaszában, vagyis az egyetemi BSc képzés keretein belül, egy kódgenerált környezetben élő préda-populáció (nyulak) viselkedésének modellezése történt meg. Az egyedek olyan örökölhető jellemzőkkel rendelkeztek, mint a mozgási sebesség, a látótávolság és különféle reprodukciós paraméterek, amelyek közvetlen hatást gyakoroltak a túlélési rátára és a természetes szelekcióra.

A egyetemi MSc képzés megkezdésével, a fejlesztés későbbi fázisába érve, a hangsúly előbb a szoftverarchitektúra optimalizálására és a refaktorálásra helyeződött. A rendszer kezdeti, monolitikus felépítését (úgynevezett „god class” probléma) egy moduláris, tiszta felelősségi körökkel rendelkező osztályszerkezet váltotta fel, jelentősen növelve a kód karbantarthatóságát és bővíthetőségét.

A biológiai hűség fokozása érdekében a szimuláció egy ragadozó fajjal (rókák) egészült ki, megteremtve a predator–prey dinamikát. A kialakuló populáció-oszcillációk vizsgálata során a Lotka–Volterra modell elméleti alapjait alkalmaztam a rendszer validálására. A szoftver legfrissebb változata már eseményvezérelt architektúrára épül, amely nemcsak a számítási hatékonyságot javítja az állapotellenőrzések minimalizálásával, hanem rugalmas keretet biztosít az új viselkedésformák implementálásához is.

A szimuláció futása során keletkező adatokat a rendszer \textbf{InfluxDB} idősoros adatbázisban tárolja, az adatok vizualizációja és a valós idejű statisztikai elemzés pedig \textbf{Grafana} segítségével valósul meg. Ez a megoldás lehetővé teszi a komplex ökológiai folyamatok transzparens nyomon követését és a különböző döntéshozatali modellek összehasonlítását.
\vfill

\selectenglish
%----------------------------------------------------------------------------
% Abstract in English
%----------------------------------------------------------------------------
\chapter*{Abstract}\addcontentsline{toc}{chapter}{Abstract}

\noindent The subject of my thesis is the development of a complex, \textbf{Unity}-based ecosystem simulation framework, primarily aimed at investigating population genetic processes and heritable traits within a digital environment. The development now spans a four-year period, during which the software has evolved from a simple agent-based model into a robust, event-driven system.

In the initial phase of the project, conducted within the framework of my BSc studies, the behavior of a prey population (rabbits) living in a code-generated environment was modeled. Individuals possessed heritable characteristics such as movement speed, vision range, and various reproductive parameters, which had a direct impact on survival rates and natural selection.

Upon commencing my MSc studies and entering a later phase of development, the focus shifted toward software architecture optimization and refactoring. The system's initial monolithic structure (the so-called "god class" problem) was replaced by a modular class hierarchy with clean responsibilities, significantly increasing code maintainability and extensibility.

To enhance biological fidelity, the simulation was expanded to include a predator species (foxes), establishing predator–prey dynamics. During the investigation of the resulting population oscillations, I applied the theoretical foundations of the Lotka–Volterra model to validate the system. The latest version of the software is built on an event-driven architecture, which not only improves computational efficiency by minimizing state checks but also provides a flexible framework for implementing new behavioral forms.

Data generated during the simulation is stored in an \textbf{InfluxDB} time-series database, while data visualization and real-time statistical analysis are realized using \textbf{Grafana}. This solution enables the transparent monitoring of complex ecological processes and the comparison of different decision-making models.
\vfill
\selectthesislanguage

\newcounter{romanPage}
\setcounter{romanPage}{\value{page}}
\stepcounter{romanPage}